% ============================================================================
% TOROIDAL EIGENMODE NAVIGATION
% ============================================================================
\subsection{Chord Co-occurrence and Markov Transition Matrix}
\label{sec:eigenmode_cooccurrence}

For each Fibonacci window scale $w \in \mathcal{W} = \{3, 5, 8, 13, 21\}$,
we build a $256 \times 256$ chord-bin co-occurrence matrix $C^{(w)}$ from
the ingested chord sequence $(\mathbf{c}_1, \ldots, \mathbf{c}_N)$:

\begin{equation}
  C^{(w)}_{ij} = \bigl|\{t : \operatorname{ChordBin}(\mathbf{c}_t) = i
    \;\land\; \exists\, t < s \leq t+w
    \text{ with } \operatorname{ChordBin}(\mathbf{c}_s) = j\}\bigr|
  \label{eq:cooccurrence}
\end{equation}

Each row is L1-normalized to form a row-stochastic Markov transition matrix:

\begin{equation}
  T^{(w)}_{ij} = \frac{C^{(w)}_{ij}}{\sum_{k=0}^{255} C^{(w)}_{ik}}
  \label{eq:markov}
\end{equation}

\subsection{Eigenvector Extraction and Phase Assignment}
\label{sec:eigen_extraction}

The eigendecomposition of $T^{(w)}$ yields eigenvalues
$\lambda_0 \geq |\lambda_1| \geq |\lambda_2| \geq \cdots$ sorted by
modulus.  The dominant eigenvalue $\lambda_0 = 1$  corresponds to the
stationary distribution and is discarded.  The two leading
\emph{subdominant} eigenvector pairs define orthogonal phase planes:

\begin{itemize}[nosep]
  \item \textbf{Primary plane (Theta):} Eigenvectors at $\lambda_1, \lambda_2$
        define $\mathbf{v}^{(w)}_{T,1}, \mathbf{v}^{(w)}_{T,2} \in \mathbb{R}^{256}$.
  \item \textbf{Secondary plane (Phi):} Eigenvectors at $\lambda_3, \lambda_4$
        define $\mathbf{v}^{(w)}_{P,1}, \mathbf{v}^{(w)}_{P,2} \in \mathbb{R}^{256}$.
\end{itemize}

If the eigenvalues form a complex conjugate pair $\lambda = \alpha \pm i\beta$,
we extract phases directly from the complex eigenvector; otherwise, we use
the two real eigenvectors as a basis.  For each structural bin $b$:

\begin{equation}
  \theta^{(w)}_b = \operatorname{atan2}\!\bigl(v^{(w)}_{T,2}[b],\;
                                               v^{(w)}_{T,1}[b]\bigr),
  \qquad
  \phi^{(w)}_b = \operatorname{atan2}\!\bigl(v^{(w)}_{P,2}[b],\;
                                              v^{(w)}_{P,1}[b]\bigr)
  \label{eq:phase_assignment}
\end{equation}

\paragraph{Eigenvector Magnitude as Frequency.}
The magnitude of the eigenvector projection measures the
\emph{structural informativeness} of each bin at scale $w$:

\begin{equation}
  f^{(w)}_{\theta,b} = 1 + \frac{\sqrt{v_{T,1}[b]^2 + v_{T,2}[b]^2}}
                                 {\max_{b'} \sqrt{v_{T,1}[b']^2 + v_{T,2}[b']^2}}
                        \cdot \kappa
  \label{eq:frequency}
\end{equation}

where $\kappa$ is the \texttt{FrequencySpread} hyperparameter controlling
the range of frequency modulation.

\subsection{Multi-Scale Circular Averaging}
\label{sec:multi_scale}

Phases from all window scales are fused via inverse-window-weighted
\emph{circular averaging} to produce scale-invariant phase assignments.
The weight for scale $w$ is:

\begin{equation}
  \alpha_w = \frac{1/w}{\sum_{w' \in \mathcal{W}} 1/w'}
  \label{eq:fib_weight}
\end{equation}

The circular mean uses the standard directional statistics formula to
avoid wraparound artifacts at $\pm\pi$:

\begin{equation}
  \Theta_b = \operatorname{atan2}\!\Bigl(
    \sum_{w \in \mathcal{W}} \alpha_w \sin\theta^{(w)}_b,\;\;
    \sum_{w \in \mathcal{W}} \alpha_w \cos\theta^{(w)}_b
  \Bigr)
  \label{eq:circular_mean}
\end{equation}

\noindent and identically for $\Phi_b$.  The final phase map
$\{\Theta_b, \Phi_b\}_{b=0}^{255}$ embeds all 256 structural bins onto
the 2-torus $\mathbb{T}^2 = S^1 \times S^1$.

\subsection{Sequential Phase Tracking and Intent Alignment}
\label{sec:toroidal_steering}

During generation, the accumulated sequential phase is tracked as a weighted
circular mean over the emitted chord sequence $(\mathbf{c}_1, \ldots,
\mathbf{c}_n)$:

\begin{equation}
  \bar{\theta}_n = \operatorname{atan2}\!\Bigl(
    \sum_{i=1}^{n} f_{\theta, b_i} \sin \Theta_{b_i},\;\;
    \sum_{i=1}^{n} f_{\theta, b_i} \cos \Theta_{b_i}
  \Bigr)
  \label{eq:seq_phase}
\end{equation}

where $b_i = \operatorname{ChordBin}(\mathbf{c}_i)$ and $f_{\theta, b_i}$
is the frequency weight.  The \emph{concentration parameter}
$\kappa_n$ (analogous to the $\kappa$ parameter of the von Mises
distribution) measures phase coherence:

\begin{equation}
  \kappa_n = \frac{\sqrt{S_n^2 + C_n^2}}{W_n}, \qquad
  S_n = \sum_{i=1}^{n} f_i \sin \Theta_{b_i}, \quad
  C_n = \sum_{i=1}^{n} f_i \cos \Theta_{b_i}, \quad
  W_n = \sum_{i=1}^{n} f_i
  \label{eq:concentration}
\end{equation}

Values of $\kappa_n$ near 1 indicate strong phase alignment (the trajectory
is coherent); values near 0 indicate phase diffusion (semantic drift).

\subsection{Geometric Closure Detection}
\label{sec:geometric_closure}

Generation terminates when the candidate chord sequence's weighted circular
mean phase returns to within a threshold $\delta$ of the anchor phase
$\bar{\theta}_{\text{anchor}}$ established at prompt initialization:

\begin{equation}
  \Delta\theta = \min\bigl(|\bar{\theta}_n - \bar{\theta}_{\text{anchor}}|,\;\;
                           2\pi - |\bar{\theta}_n - \bar{\theta}_{\text{anchor}}|\bigr)
  < \delta
  \label{eq:closure}
\end{equation}

with $\delta = 0.45$ radians.  This constitutes a topological proof of
structural closure: the generative trajectory has completed a semantically
coherent loop on the torus.  Because $\mathbb{T}^2$ is a compact manifold,
every trajectory is pre-compact (bounded), and the density of the
trajectory---guaranteed by the irrationality of the eigenvalue ratios for
generic data---prevents entrapment in isolated fixed points or periodic
orbits.
